\section{Opus Deployment Prompt --- §2.6 US Profitability Analysis:
Final
Pass}\label{opus-deployment-prompt--26-us-profitability-analysis-final-pass}

\subsection{Purpose: Parsimonious rewrite with full voice
calibration}\label{purpose-parsimonious-rewrite-with-full-voice-calibration}

\subsection{Date: 2026-04-07}\label{date-2026-04-07}

\subsection{Instruction: Load everything below before writing a single
word of
output.}\label{instruction-load-everything-below-before-writing-a-single-word-of-output}

\begin{center}\rule{0.5\linewidth}{0.5pt}\end{center}

\subsection{YOUR TASK}\label{your-task}

You are rewriting §2.6.2 of a heterodox political economy dissertation
chapter. The input draft (v2, \textasciitilde3,100 words) is provided in
full at the end of this prompt. Your output should be approximately
\textbf{1,600--1,800 words} of final-quality LaTeX prose --- a reduction
of roughly 45 percent --- achieved entirely through compression and
voice calibration, not through removing empirical content.

The rewrite has three simultaneous objectives:

\begin{enumerate}
\def\labelenumi{\arabic{enumi}.}
\tightlist
\item
  \textbf{Compress} without losing any empirical finding or theoretical
  claim
\item
  \textbf{Calibrate voice} to match the author\textquotesingle s
  linguistic and rhetorical patterns
\item
  \textbf{Foreground the regulationist framing} of functionality and
  dysfunctionality as the organizing spine --- this is the theoretical
  contribution that v2 states but does not sufficiently load into the
  prose
\end{enumerate}

\begin{center}\rule{0.5\linewidth}{0.5pt}\end{center}

\subsection{PART 1 --- STYLE CALIBRATION: WLM
3.1}\label{part-1--style-calibration-wlm-31}

The author\textquotesingle s calibrated writing style vector operates in
a 16-dimensional space. The scores below govern how you write. Protected
strengths (≥9) must not be diluted. Revision targets (\textless8) should
be actively improved.

\subsubsection{Protected strengths --- never
weaken}\label{protected-strengths--never-weaken}

\begin{longtable}[]{@{}llll@{}}
\toprule\noalign{}
Code & Dimension & Score & What it means in practice \\
\midrule\noalign{}
\endhead
\bottomrule\noalign{}
\endlastfoot
TP & Theoretical Positioning Precision & 9.6 & Every empirical claim
sits inside a theoretical category; numbers float toward structural
interpretations \\
EA & Empirical Admissibility Discipline & 9.5 & No causal language; no
inferential overclaiming; "records," "registers," "accounts for" ---
never "drives," "shows," "proves" \\
TS2 & Terminology Stability & 9.6 & Terms locked in the jargon table
below; no synonyms introduced mid-section \\
CD & Concept Definition Discipline & 9.3 & Terms carry their meaning on
first use; never introduced by symbol alone \\
FI & Formal-Interpretive Bridge Quality & 9.2 & Every number connects
immediately to a structural category; no number floats \\
\end{longtable}

\subsubsection{Revision targets --- actively
improve}\label{revision-targets--actively-improve}

\begin{longtable}[]{@{}llll@{}}
\toprule\noalign{}
Code & Dimension & Score & What improvement looks like \\
\midrule\noalign{}
\endhead
\bottomrule\noalign{}
\endlastfoot
SG & Signposting Quality & 7.8 & Each paragraph opens with its
structural claim stated as a navigational sentence \\
CH & Citation Hygiene & 7.9 & Citations mid-sentence attached to
specific claims; never end-clustered \\
CE & Compression Efficiency & 7.9 & No sentence restates what adjacent
sentence already said; reference tables once \\
RR & Redundancy Control & 8.0 & One claim per paragraph, developed
through mechanism → evidence → consequence \\
\end{longtable}

\begin{center}\rule{0.5\linewidth}{0.5pt}\end{center}

\subsection{PART 2 --- AUTHOR VOICE: LINGUISTIC AND RHETORICAL
PATTERNS}\label{part-2--author-voice-linguistic-and-rhetorical-patterns}

These are the structural signatures of the author\textquotesingle s
prose, extracted from Chapter 1
(\texttt{Chapter1\_CriticalReplication.pdf}). Reproduce them.

\subsubsection{Sentence architecture}\label{sentence-architecture}

The author writes long, architecturally complex sentences with multiple
embedded subordinate clauses that carry their own evidence without
becoming run-ons. The subordinate clauses do the explanatory work; the
main clause states the structural claim.

Example from Chapter 1:

\begin{quote}
"The FRB series was built to produce a stable, institutionally credible
signal for monetary policy governance within a historically specific
Fordist settlement that made slack legible, actionable, and
contestable."
\end{quote}

The sentence structure is: {[}structural claim{]} within
{[}regime-specific qualifier{]} that {[}consequence{]}. Not:
{[}claim{]}. {[}Explanation{]}. {[}Consequence{]}. Three separate
sentences would dilute the density. One complex sentence preserves it.

\subsubsection{The em-dash as elaborative
extension}\label{the-em-dash-as-elaborative-extension}

Em-dashes extend claims without interrupting them. They introduce a
parenthetical that adds precision or mechanism to what precedes it, not
a detour from it.

Example from Chapter 1:

\begin{quote}
"What emerged, then, was not simply a better estimate of productive
ceilings, but a routinized operational device through which firms,
analysts, and policymakers could speak a common language about slack."
\end{quote}

\subsubsection{The "not X but Y" precision
construction}\label{the-not-x-but-y-precision-construction}

The author uses negative-first constructions to specify exactly what
something is by first ruling out what it is not:

\begin{quote}
"The Fordist settlement did not abolish that routine. It triggered the
process of its recodification."
\end{quote}

Apply this to empirical findings: "The 1972--74 contraction does not
record a structural deterioration in productive efficiency --- it
records a price-mediated compression of the nominal rate of return."

\subsubsection{Active verbs drawn from the regulationist
vocabulary}\label{active-verbs-drawn-from-the-regulationist-vocabulary}

Prefer: \textbf{registered, routinized, sedimented, embedded,
institutionalized, displaced, compressed, exhausted, contained,
reconfigured, accelerated, activated, generated, sustained, eroded}.

Avoid: showed, demonstrated, revealed, proved, indicated, suggested.

\subsubsection{Citation placement}\label{citation-placement}

Citations appear mid-sentence, attached to the specific claim they
support. Never cluster at paragraph end. Never introduce a paragraph
with a citation.

Correct: "...connecting the distributional channel to the institutional
ceiling of full employment itself \textbackslash citep\{Baker1996,
Vidal2019\}."

Incorrect: "The distributional channel was eroded by labour-market
tightening. See \textbackslash citet\{Baker1996\} and
\textbackslash citet\{Vidal2019\}."

\subsubsection{Paragraph closure}\label{paragraph-closure}

Final sentences draw consequences, not summaries. Never restate the
topic sentence. The closing sentence should open a logical door toward
the next paragraph without explicitly announcing that it does.

\subsubsection{Compression cadence}\label{compression-cadence}

Compression is not achieved by imposing rigid word budgets on each
movement. It is achieved by raising the density of each sentence so that
one complex sentence carries what v2 used two or three sentences for.
The author\textquotesingle s natural cadence alternates between
architecturally complex sentences (long, multi-clause, structurally
loaded) and short declarative sentences that land a claim after the
complex sentence has built the scaffolding. Vary the rhythm: do not
write seven complex sentences in a row. Let short sentences punctuate.

\begin{center}\rule{0.5\linewidth}{0.5pt}\end{center}

\subsection{PART 3 --- JARGON CALIBRATION: LOCKED
SUBSTITUTIONS}\label{part-3--jargon-calibration-locked-substitutions}

These replacements are mandatory throughout. No exceptions.
Back-references to locked terms retain their full attribution on every
use, not just first use.

\begin{longtable}[]{@{}ll@{}}
\toprule\noalign{}
Write & Never write \\
\midrule\noalign{}
\endhead
\bottomrule\noalign{}
\endlastfoot
"overaccumulation regime (θ \textless{} 1)" & "sub-Harrodian" \\
"underaccumulation regime (θ \textgreater{} 1)" & "super-Harrodian" \\
"productive capacities" (plural always) & "productive capacity"
(singular) \\
"accumulation dynamics" & "growth dynamics" / "growth trajectory" \\
"dysfunctionality of the accumulation regime" & "regime failure" /
"breakdown" \\
"stagnation tendencies in Vidal\textquotesingle s (2014) sense" &
"stagnation tendency" standalone --- every use must carry attribution \\
"crisis-trigger mechanism (Okishio et al. 2022)" & "crisis trigger"
standalone \\
"the corresponding cell of the Basu (2019) typology" & "Basu cell" \\
"institutional settlement" / "institutional configuration" &
"institutions" generic \\
"transformation elasticity" (written out on first use) & θ alone in
prose \\
"knife-edge at θ = 1" (empirical sections only) & "Harrodian benchmark"
in theory prose \\
\end{longtable}

Freely admissible: "dysfunctionality," "functionality,"
"overaccumulation," "accumulation regime," "regulationist,"
"routinized," "recodification," "institutional settlement," "stagnation
tendencies" (with Vidal attribution).

Pre-estimation register (active throughout):

\begin{itemize}
\tightlist
\item
  Past tense for what the decomposition records
\item
  Present tense for structural claims and mechanisms
\item
  "The decomposition records," "the accounting identity reveals," "the
  shares register" --- never "we find," "results show," "the data
  demonstrates"
\end{itemize}

\begin{center}\rule{0.5\linewidth}{0.5pt}\end{center}

\subsection{PART 4 --- REGULATIONIST FRAMING: THE ORGANIZING
SPINE}\label{part-4--regulationist-framing-the-organizing-spine}

This is the primary calibration instruction for the rewrite. The v2
draft states the functionality/dysfunctionality distinction but does not
make it the \textbf{organizing principle} of the prose. The rewrite
must.

\subsubsection{Core concepts}\label{core-concepts}

\textbf{Functionality of the accumulation regime} means that the
institutional settlement successfully contains the stagnation tendencies
the regime generates. In channel terms: when one channel registers a
stagnation tendency in Vidal\textquotesingle s (2014) sense (downward
pressure on profitability), the others offset it --- the net effect is a
manageable contraction that partial institutional reconfiguration can
address. Single-channel contractions are the empirical form of
stagnation tendencies within a functional regime.

\textbf{Dysfunctionality} means that multiple stagnation tendencies
reinforce simultaneously, overwhelming the containment capacity of the
institutional settlement. In channel terms: all channels drag together,
no channel offsets, and no single-channel reconfiguration suffices. The
1966--1970 three-channel reinforcing structure is the empirical form of
dysfunctionality.

\textbf{Offsetting vs. reinforcing} is the key analytical distinction:

\begin{itemize}
\tightlist
\item
  Offsetting = one tendency active, others neutral or compensating →
  functional
\item
  Reinforcing = all tendencies active simultaneously → dysfunctional
\end{itemize}

\subsubsection{How to apply this to the
draft}\label{how-to-apply-this-to-the-draft}

The prose should be organized around a single question: \textbf{was the
accumulation regime, in each sub-period, containing or failing to
contain its stagnation tendencies?}

Sub-period averages (Table B3): the profit rate\textquotesingle s
cross-sub-period stability is the macroscopic sign of a regime that
contained its stagnation tendencies --- \(B_{r}\) rising offset falling
\(\hat{\mu}\) and \(\pi\). This is functionality registered at the
sub-period level.

Early contractions (1951--58): single-channel dominance means the regime
was functional --- stagnation tendencies did not reinforce, and partial
reconfiguration was structurally sufficient.

Great Fordist Expansion (1958--66): all channels positive means the
regime was not merely containing tendencies but actively generating
profitability from all four structural channels simultaneously --- the
high-water mark of functionality.

Structural crisis (1966--70): three-channel reinforcing means the regime
became dysfunctional --- stagnation tendencies that had previously been
contained (distributional pressure, productivity deceleration, demand
compression) now reinforced simultaneously, and the institutional
settlement could not contain their combination.

Oil shock (1972--74): an external perturbation of a structurally
exhausted regime, not a stagnation tendency generated internally. The
\(Py/PK\) dominance (57.8\%) is not a product of the
regime\textquotesingle s own dynamics but of an exogenous price shock
--- which is why the oil shock accelerated institutional collapse rather
than initiating it. The distinction matters: the regime was already
dysfunctional in 1966--70; 1972--74 is its termination, not its crisis.

\subsubsection{Crisis-theory framing (from
§2.2)}\label{crisis-theory-framing-from-22}

The results should be interpreted in light of the multi-scalar crisis
theory developed earlier in the chapter. The key conceptual layers are:

\begin{enumerate}
\def\labelenumi{\arabic{enumi}.}
\item
  \textbf{Shaikh (1978) / Itoh (1978)}: Crisis is endemic to capitalism;
  the limit of accumulation is capitalism itself. Itoh\textquotesingle s
  distinction between \emph{excess of capital} (overaccumulation) and
  \emph{excess commodity} (realization failure) maps onto different
  channel configurations in the decomposition.
\item
  \textbf{Okishio (1961; 2022)}: The crisis trigger is the cumulative
  path of \(\mu_t\) --- the variable that registers the deviation of
  actual accumulation from the accumulation equilibrium trajectory. When
  the classical long-run closure is lifted, \(\mu_t\) emerges as
  structurally determined and cumulatively unstable.
\item
  \textbf{Basu (2019) typology}: Each swing can be mapped onto a cell of
  the surplus-value accounting grid (underconsumption / profit squeeze /
  financial fragility / falling rate of profit). Channel dominance
  determines the cell assignment.
\item
  \textbf{Vidal (2014, 2019)}: Stagnation tendencies, partial crisis,
  structural crisis as the three categories. The key regulationist move
  is that offsetting tendencies within a functional regime are not
  equilibrium --- they are the institutional containment of
  contradictions that remains operative only as long as the settlement
  holds.
\item
  \textbf{Vidal (2022) / firm-level}: The Taylorist labour process that
  anchored the Fordist settlement contained its own dissolution --- the
  management contradiction (coordination vs. discipline) and workforce
  contradiction (productive socialisation vs. alienation) meant that the
  same deskilling that served accumulation in the extensive phase became
  a fetter in the intensive phase. This is the micro-foundation of the
  \(B_r\) channel\textquotesingle s exhaustion.
\end{enumerate}

When the prose describes "full employment" or "near-full employment,"
this is admissible as historical description grounded in the crisis
theory: full employment names the institutional ceiling at which the
wage share approaches its feasible maximum
(\(\omega \to \omega_{\max}\)), the working class\textquotesingle s
bargaining leverage intensifies, and the distributional channel comes
under structural pressure. The formal analytical content
(\(\omega \to \omega_{\max}\)) and the historiographical description
("full-employment conditions") serve each other --- use whichever is
more precise in context, or both when one grounds the other.

\subsubsection{Stage boundary: θ and the distributional
variable}\label{stage-boundary-ux3b8-and-the-distributional-variable}

In this section, the transformation elasticity θ has already been
estimated in Stage A and was the mediator through which \(\hat{\mu}_t\)
was structurally identified. It is not discussed further here --- Stage
A is finished. Do not re-introduce the ω-endogeneity of θ or any Stage A
identification issues.

The distributional variable \(\pi\) (profit share) used throughout this
section is the Weisskopf decomposition channel --- an accounting object.
It is the reciprocal of the wage share ω used in Stage A
(\(\pi + \omega = 1\)). The identification in Stage A used ω because the
cointegrating vector and its coefficient dynamics are more naturally
interpreted in wage-share space. Here we use \(\pi\) because it is the
accounting variable from the Weisskopf decomposition, and it was decided
not to re-frame it as wage share for the profitability analysis.
Understand and respect this boundary between stages: do not flag \(\pi\)
as inconsistent with the Stage A specification.

\begin{center}\rule{0.5\linewidth}{0.5pt}\end{center}

\subsection{PART 5 --- COMPRESSION
STRATEGY}\label{part-5--compression-strategy}

Target: 1,600--1,800 words from \textasciitilde3,100. Compression comes
from density --- one sentence carrying what v2 used two sentences for
--- not from removal of findings. Do not impose rigid per-movement word
budgets. Instead, let voice calibration and natural cadence determine
where prose expands and where it compresses. The following are
structural priorities, not word counts:

\begin{itemize}
\tightlist
\item
  \textbf{Opening (methodology)}: Maximum compression. 2 sentences.
  Framework assumed from §2.3.
\item
  \textbf{A (structural background)}: Merge A1+A2 into a single
  paragraph. Keep compensation mechanism + PyPK headwind as the two
  claims. This is setup, not the finding.
\item
  \textbf{B (turning-point structure)}: Brief. What B4 shows, why r\_t
  not GDP, the two-phase structure. 3--4 sentences.
\item
  \textbf{C1+C2 (functional regime)}: Merge into one paragraph. Early
  contractions as single-channel → functionality; expansion as coherence
  peak. Give enough room for the demand-led mechanism of 1958--66 to
  land.
\item
  \textbf{C3 (dysfunctionality)}: This is the central finding. Give it
  the most room. The three-mechanism convergence (distributional
  pressure, Taylorist exhaustion, demand ceiling) must be stated with
  full force.
\item
  \textbf{C4 (partial recovery + reattribution)}: Partial recovery
  compressed to 2 sentences. The 1972--74 reattribution is the
  chapter\textquotesingle s technical contribution at Stage B --- give
  it enough room to land clearly.
\item
  \textbf{D (synthesis)}: Channel progression + compensation
  contradiction + punchline. One paragraph.
\item
  \textbf{Transition}: 2--3 sentences bridging to Stage C.
\end{itemize}

The "Movement" labels in the v2 draft comments are for your orientation
only --- they do not appear in the output.

\begin{center}\rule{0.5\linewidth}{0.5pt}\end{center}

\subsection{PART 6 --- EMPIRICAL
ANCHORS}\label{part-6--empirical-anchors}

These numbers and claims must survive compression exactly as stated. Do
not paraphrase, round, or omit.

\textbf{Note on stage boundary:} \(\pi\) here is the profit share
\emph{channel} in the Weisskopf decomposition, not the distributional
state variable from Stage A. They are reciprocal (\(\pi + \omega = 1\)).
Stage A identified in \(\omega\)-space for coefficient interpretation;
Stage B uses \(\pi\) because it is the accounting variable from the
decomposition.

From Table B3 (sub-period means):

\begin{itemize}
\tightlist
\item
  r: 0.187 (Early Fordist) / 0.189 (Late Fordist) --- stable
\item
  μ: 0.737 → 0.723 → 0.686 (monotonic decline)
\item
  B\_real: 0.586 → 0.635 (structural rise)
\item
  PyPK: 1.369 → 1.176 → 1.140 → 1.025 (monotonic decline, 25pp total)
\item
  π: 0.369 → 0.360 (modest compression)
\end{itemize}

From Table B2 (swing shares, 4-channel):

\begin{itemize}
\tightlist
\item
  1951--54 contraction: π dominant (48.5\%), μ secondary (41.7\%), no
  offsets
\item
  1954--55 expansion: B\_real dominant (77.7\%), PyPK offsets (−47.2\%)
\item
  1955--58 contraction: μ dominant (60.2\%), no offsets
\item
  1958--66 expansion: μ dominant (47.7\%), B\_real (22.4\%), PyPK
  (13.2\%), π (16.7\%), no offsets
\item
  1966--70 contraction: π (34.5\%), B\_real (29.2\%), μ (29.8\%), all
  drag, no offsets --- THREE-WAY
\item
  1970--72 expansion: π dominant (52.7\%), PyPK offsets (−10.4\%)
\item
  1972--74 contraction: PyPK dominant (57.8\%), B\_real offsets
  (−13.0\%) --- OIL SHOCK REATTRIBUTION
\end{itemize}

The 1972--74 reattribution is the chapter\textquotesingle s technical
contribution at Stage B. It must be stated clearly: what three-channel
decompositions attribute to technology, the four-channel decomposition
attributes to the relative price channel. Real capital productivity
(B\_real) offset during the oil shock contraction --- it did not
collapse.

\begin{center}\rule{0.5\linewidth}{0.5pt}\end{center}

\subsection{PART 7 --- EXECUTION RULES}\label{part-7--execution-rules}

\begin{enumerate}
\def\labelenumi{\arabic{enumi}.}
\item
  Write in LaTeX. Use the same environments as the input draft. Preserve
  all \texttt{\textbackslash{}label\{\}},
  \texttt{\textbackslash{}ref\{\}}, and
  \texttt{\textbackslash{}citet\{\}}/\texttt{\textbackslash{}citep\{\}}
  commands.
\item
  Keep the shared source note
  \texttt{\textbackslash{}begin\{quote\}...\textbackslash{}end\{quote\}}
  exactly as in v2.
\item
  Do not add sections, sub-sections, or
  \texttt{\textbackslash{}paragraph\{\}} commands not in v2. The section
  structure is: opening → A → B → C1+C2 (merged) → C3 → C4 → D →
  transition.
\item
  Do not use bullet points anywhere in the body prose.
\item
  The closing sentence of the synthesis (Movement D) must be exactly:
  "Demand conditions intervened on the political economy of the ceiling
  -\/-\/- the class struggle over distribution, utilisation, and the
  institutional forms that govern whether accumulation can proceed."
  This sentence is locked. Do not paraphrase it.
\item
  \textbf{Correct any surviving jargon violations from the table in Part
  3 during the rewrite.} In particular: "productive capacity" (singular)
  → "productive capacities" (plural always); "stagnation tendency"
  without Vidal attribution → add attribution.
\item
  \textbf{Back-references to locked terms retain their full attribution
  on every use, not just first use.} Do not drop "in
  Vidal\textquotesingle s (2014) sense" after the first mention.
\item
  \textbf{"Movement" labels in the v2 draft\textquotesingle s LaTeX
  comments are for your orientation.} They do not appear in the output.
  Do not produce \texttt{\%\ MOVEMENT} comment blocks in the output ---
  write clean LaTeX with minimal structural comments.
\item
  After writing the draft, append a brief self-audit (in a comment block
  \texttt{\%\ AUDIT:\ ...}) noting: (a) actual word count, (b) any
  compression decisions that sacrificed a specific claim, (c) any jargon
  violation flagged and resolved.
\end{enumerate}

\begin{center}\rule{0.5\linewidth}{0.5pt}\end{center}

\subsection{INPUT DRAFT (v2 --- full text for
rewriting)}\label{input-draft-v2--full-text-for-rewriting}

\% s26\_us\_draft\_v2.tex \% §2.6 Stage B: Profitability Analysis ---
United States \% Second pass: WLM 3.1 + Jargon Calibration +
hetpe-academic-writing \% Explanatory mode on --- framework assumed
established in §2.3 \% Date: 2026-04-06

\textbackslash subsubsection\{Stage B: Profitability Analysis --- United
States\} \textbackslash label\{sssec:stageB\_us\}

\%
-\/-\/-\/-\/-\/-\/-\/-\/-\/-\/-\/-\/-\/-\/-\/-\/-\/-\/-\/-\/-\/-\/-\/-\/-\/-\/-\/-\/-\/-\/-\/-\/-\/-\/-\/-\/-\/-\/-\/-\/-\/-\/-\/-\/-\/-\/-\/-\/-\/-\/-\/-\/-\/-\/-\/-\/-\/-\/-\/-\/-\/-\/-\/-\/-\/-\/-\/-\/-\/-\/-
\% OPENING --- framework assumed; methodology compressed to 3 sentences
\%
-\/-\/-\/-\/-\/-\/-\/-\/-\/-\/-\/-\/-\/-\/-\/-\/-\/-\/-\/-\/-\/-\/-\/-\/-\/-\/-\/-\/-\/-\/-\/-\/-\/-\/-\/-\/-\/-\/-\/-\/-\/-\/-\/-\/-\/-\/-\/-\/-\/-\/-\/-\/-\/-\/-\/-\/-\/-\/-\/-\/-\/-\/-\/-\/-\/-\/-\/-\/-\/-\/-

The four-channel decomposition applied here extends the Weisskopf (1979)
framework developed in
Section\textasciitilde\textbackslash ref\{sec:framework\} by separating
nominal capital productivity into its relative price and real
components, so that the exact identity
\$\textbackslash Delta\textbackslash ln r =
\textbackslash phi\^{}\{\textbackslash mu\} +
\textbackslash phi\^{}\{Py/PK\} + \textbackslash phi\^{}\{B\_\{r\}\}

\begin{itemize}
\tightlist
\item
  \textbackslash phi\^{}\{\textbackslash pi\}\$ replaces the standard
  three-channel form. The capacity utilization series \(\hat{\mu}_{t}\)
  is the structurally identified estimate from Stage\textasciitilde A
  rather than a filter residual, which is what permits the demand and
  technology channels to be read as analytically distinct objects rather
  than composites of the same identification artefact. The note on
  figures preceding this sub-section records the measurement
  conventions; what follows is their interpretation.
\end{itemize}

\% Note on figures \textbackslash begin\{quote\}\textbackslash small
\textbackslash textit\{Note on figures:\} All figures cover the US
Non-Financial Corporate Sector, 1940-\/-1978.
\(r_{t} = \text{GOS}_{t}/\text{KNC}_{t}\); \(\hat{\mu}_{t}\) from
Stage\textasciitilde A MPF (pinned at \(\hat{\mu}_{1948} = 0.80\));
\(B_{r,t} = \text{GVA}_{r,t}/(\text{KNR}_{t}\cdot\hat{\mu}_{t})\);
\(Py/PK_{t}\) = GDP deflator / capital price deflator. Regime shading:
Pre-Fordist (1940-\/-1946), Fordist (1947-\/-1973), Post-Fordist onset
(1974-\/-1978). Period windows: WWII buildup (1940-\/-1944), Interim
(1944-\/-1947). \textbackslash end\{quote\}

\%
-\/-\/-\/-\/-\/-\/-\/-\/-\/-\/-\/-\/-\/-\/-\/-\/-\/-\/-\/-\/-\/-\/-\/-\/-\/-\/-\/-\/-\/-\/-\/-\/-\/-\/-\/-\/-\/-\/-\/-\/-\/-\/-\/-\/-\/-\/-\/-\/-\/-\/-\/-\/-\/-\/-\/-\/-\/-\/-\/-\/-\/-\/-\/-\/-\/-\/-\/-\/-\/-\/-
\% MOVEMENT A1 --- THE COMPENSATION MECHANISM \%
-\/-\/-\/-\/-\/-\/-\/-\/-\/-\/-\/-\/-\/-\/-\/-\/-\/-\/-\/-\/-\/-\/-\/-\/-\/-\/-\/-\/-\/-\/-\/-\/-\/-\/-\/-\/-\/-\/-\/-\/-\/-\/-\/-\/-\/-\/-\/-\/-\/-\/-\/-\/-\/-\/-\/-\/-\/-\/-\/-\/-\/-\/-\/-\/-\/-\/-\/-\/-\/-\/-

The sub-period means in
Table\textasciitilde\textbackslash ref\{tab:B3\_subperiod\} reveal a
puzzle that goes to the core of the Fordist political economy: the
nominal profit rate was essentially flat across the two Fordist
sub-periods -\/-\/- 0.187 in the Early Fordist (1945-\/-1966) and 0.189
in the Late Fordist (1967-\/-1973) -\/-\/- even as two of its three
constituent channels were contracting. Mean capacity utilization fell
from 0.737 to 0.723 across the same sub-period boundary, registering the
secular deterioration in the activation of productive capacities that
Stage\textasciitilde A\textquotesingle s overaccumulation finding
predicts: with the transformation elasticity linking capital
accumulation to productive capacity expansion consistently below unity,
each successive wave of investment expanded the productive frontier more
slowly than it expanded the capital stock, widening the structural gap
between what the economy could produce and what it actually realized.
Mean profit share fell from 0.369 to 0.360, recording the distributional
pressure that the political economy of full employment imposed as the
Fordist settlement reached its later phase. The channel that offset both
declines was real capital productivity at normal capacity, which rose
from 0.586 to 0.635 across the same transition. What this means is that
the Fordist accumulation regime reproduced profitability not by
sustaining demand activation or protecting the profit share but by
continuously raising the real productive efficiency with which each unit
of capital stock generated output at normal utilization -\/-\/- the
technology channel doing the compensatory work that the demand and
distributional channels could no longer perform.
Figures\textasciitilde\textbackslash ref\{fig:B1b\} and
\textbackslash ref\{fig:B1c\} make this offsetting trajectory visible in
indexed form: the \(\hat{\mu}\) series declines from its wartime peak
toward the Fordist plateau and continues falling through the 1970s,
while the \(B_{r}\) series climbs steadily through the entire Fordist
window, the two trajectories moving in opposite directions around the
stable profit rate that their product sustains.

\%
-\/-\/-\/-\/-\/-\/-\/-\/-\/-\/-\/-\/-\/-\/-\/-\/-\/-\/-\/-\/-\/-\/-\/-\/-\/-\/-\/-\/-\/-\/-\/-\/-\/-\/-\/-\/-\/-\/-\/-\/-\/-\/-\/-\/-\/-\/-\/-\/-\/-\/-\/-\/-\/-\/-\/-\/-\/-\/-\/-\/-\/-\/-\/-\/-\/-\/-\/-\/-\/-\/-
\% MOVEMENT A2 --- THE RELATIVE PRICE HEADWIND \%
-\/-\/-\/-\/-\/-\/-\/-\/-\/-\/-\/-\/-\/-\/-\/-\/-\/-\/-\/-\/-\/-\/-\/-\/-\/-\/-\/-\/-\/-\/-\/-\/-\/-\/-\/-\/-\/-\/-\/-\/-\/-\/-\/-\/-\/-\/-\/-\/-\/-\/-\/-\/-\/-\/-\/-\/-\/-\/-\/-\/-\/-\/-\/-\/-\/-\/-\/-\/-\/-\/-

Reading the compensation mechanism through quantity channels alone,
however, misses a structural feature that the four-channel decomposition
makes visible and the three-channel framework suppresses: the
output-to-capital relative price was compressing monotonically
throughout the Fordist period, imposing a price-denominated headwind
against profitability that the real capital productivity channel had to
overcome before it could begin to compensate for falling demand
activation. Table\textasciitilde\textbackslash ref\{tab:B3\_subperiod\}
records mean \(Py/PK\) falling from 1.369 in the Pre-Fordist/WWII period
to 1.176 in the Early Fordist, 1.140 in the Late Fordist, and 1.025 at
the Post-Fordist onset -\/-\/- a decline of 25 percentage points across
the full window. The implication is that capital goods prices were
rising faster than output prices as a structural feature of Fordist
accumulation dynamics, which means that the nominal rate of return on
each additional unit of capital was being squeezed from the price side
even before the oil shock of 1972-\/-74 made that compression
dramatically visible. The oil shock did not initiate this tendency; it
accelerated a pre-existing secular movement to the point where the price
channel displaced the technology channel as the dominant driver of the
profit rate contraction, a finding to which the swing-level
decomposition will return in Paragraph\textasciitilde C4. For now, the
point is that the compensation mechanism documented in
Paragraph\textasciitilde A1 operated within a price environment that was
itself working against profitability recovery -\/-\/- meaning that
\(B_{r}\) had to grow not merely to offset declining \(\hat{\mu}\) but
to do so while fighting a structural relative price deterioration that
the three-channel nominal framework attributes entirely to the
technology channel without distinguishing these two qualitatively
different sources of capital productivity pressure.

\%
-\/-\/-\/-\/-\/-\/-\/-\/-\/-\/-\/-\/-\/-\/-\/-\/-\/-\/-\/-\/-\/-\/-\/-\/-\/-\/-\/-\/-\/-\/-\/-\/-\/-\/-\/-\/-\/-\/-\/-\/-\/-\/-\/-\/-\/-\/-\/-\/-\/-\/-\/-\/-\/-\/-\/-\/-\/-\/-\/-\/-\/-\/-\/-\/-\/-\/-\/-\/-\/-\/-
\% MOVEMENT B --- TURNING-POINT STRUCTURE \%
-\/-\/-\/-\/-\/-\/-\/-\/-\/-\/-\/-\/-\/-\/-\/-\/-\/-\/-\/-\/-\/-\/-\/-\/-\/-\/-\/-\/-\/-\/-\/-\/-\/-\/-\/-\/-\/-\/-\/-\/-\/-\/-\/-\/-\/-\/-\/-\/-\/-\/-\/-\/-\/-\/-\/-\/-\/-\/-\/-\/-\/-\/-\/-\/-\/-\/-\/-\/-\/-\/-

Against this structural background,
Figure\textasciitilde\textbackslash ref\{fig:B4\} identifies seven
turning points on the nominal profit rate between 1951 and 1974 that
organize the Fordist accumulation dynamics into the swing-level episodes
whose channel composition
Table\textasciitilde\textbackslash ref\{tab:B2\_peaktotrough\}
decomposes. Turning points are identified directly on \(r_{t}\) rather
than on output or employment because the object of the decomposition is
the profit rate, not the business cycle, and the two do not always share
their timing -\/-\/- a profit rate contraction can begin while output is
still expanding if the capital stock is growing faster than output,
which is precisely the structural condition the overaccumulation regime
generates. The seven-swing sequence spans two qualitatively distinct
structural phases. The first four swings (1951-\/-1966) constitute the
Fordist settlement\textquotesingle s normal cycle: moderate-amplitude
contractions followed by recoveries that cumulatively lifted the profit
rate to its 1966 peak. The remaining three swings (1966-\/-1974)
constitute its unravelling: a deep contraction that the partial recovery
of 1970-\/-1972 could not undo, terminated by an external shock that
compressed what little room for manoeuvre the institutional settlement
retained. The largest expansion in the sample (\(\Delta\ln r = +0.336\),
1958-\/-1966) and the most consequential contraction
(\(\Delta\ln r = -0.269\), 1966-\/-1970) define the structural arc;
everything else in the decomposition is a commentary on why those two
swings had the channel structure they did.

\%
-\/-\/-\/-\/-\/-\/-\/-\/-\/-\/-\/-\/-\/-\/-\/-\/-\/-\/-\/-\/-\/-\/-\/-\/-\/-\/-\/-\/-\/-\/-\/-\/-\/-\/-\/-\/-\/-\/-\/-\/-\/-\/-\/-\/-\/-\/-\/-\/-\/-\/-\/-\/-\/-\/-\/-\/-\/-\/-\/-\/-\/-\/-\/-\/-\/-\/-\/-\/-\/-\/-
\% MOVEMENT C1 --- THE EARLY CONTRACTIONS \%
-\/-\/-\/-\/-\/-\/-\/-\/-\/-\/-\/-\/-\/-\/-\/-\/-\/-\/-\/-\/-\/-\/-\/-\/-\/-\/-\/-\/-\/-\/-\/-\/-\/-\/-\/-\/-\/-\/-\/-\/-\/-\/-\/-\/-\/-\/-\/-\/-\/-\/-\/-\/-\/-\/-\/-\/-\/-\/-\/-\/-\/-\/-\/-\/-\/-\/-\/-\/-\/-\/-

The 1951-\/-1954 and 1955-\/-1958 contractions record single-channel
dominance, which is the analytical signature of stagnation tendencies in
Vidal\textquotesingle s (2014) sense -\/-\/- downward pressures on
profitability that a functional accumulation regime can absorb through
partial institutional reconfiguration precisely because only one
structural channel is under serious pressure at a time. In the
1951-\/-1954 contraction (\(\Delta\ln r
= -0.133\)), the distributional channel accounts for 48.5 percent of the
compounded decline and the demand channel for 41.7 percent, with real
capital productivity (5.7 percent) and the relative price channel (4.1
percent) negligible. The distributional dominance places this episode in
the profit squeeze cell of the Basu (2019) typology: Korean War wage
pressure tightened labour markets and compressed the profit share while
military spending wound down, and the two leading channels dragged
simultaneously without any offsetting force. What makes this a
stagnation tendency rather than a structural crisis, in
Vidal\textquotesingle s taxonomy, is that a distributional adjustment
-\/-\/- a restoration of wage restraint as the wartime labour-market
conditions eased -\/-\/- was structurally sufficient to address the
dominant pressure. The 1955-\/-1958 contraction (\(\Delta\ln r
= -0.191\)) records a different channel structure but the same
single-channel logic: the demand channel accounts for 60.2 percent of
the compounded decline, with distribution contributing 25.3 percent and
both the technology and relative price channels negligible. This is the
crisis-trigger mechanism identified by
\textbackslash citet\{OkishioEtAl2022\} most clearly operative in the
Fordist sample -\/-\/- productive capacities were in place but could not
be activated because effective demand contracted, and the realization
failure is registered in the profit rate through the \(\hat{\mu}\)
channel. The single-channel dominance again signals resolvability:
demand stimulus -\/-\/- which the Eisenhower recession eventually
received through federal spending adjustments -\/-\/- could in principle
address a \(\mu\)-dominant contraction without restructuring the
institutional settlement. Both episodes are, in this precise sense, the
normal pathology of a functional Fordist accumulation regime.

\%
-\/-\/-\/-\/-\/-\/-\/-\/-\/-\/-\/-\/-\/-\/-\/-\/-\/-\/-\/-\/-\/-\/-\/-\/-\/-\/-\/-\/-\/-\/-\/-\/-\/-\/-\/-\/-\/-\/-\/-\/-\/-\/-\/-\/-\/-\/-\/-\/-\/-\/-\/-\/-\/-\/-\/-\/-\/-\/-\/-\/-\/-\/-\/-\/-\/-\/-\/-\/-\/-\/-
\% MOVEMENT C2 --- THE GREAT FORDIST EXPANSION \%
-\/-\/-\/-\/-\/-\/-\/-\/-\/-\/-\/-\/-\/-\/-\/-\/-\/-\/-\/-\/-\/-\/-\/-\/-\/-\/-\/-\/-\/-\/-\/-\/-\/-\/-\/-\/-\/-\/-\/-\/-\/-\/-\/-\/-\/-\/-\/-\/-\/-\/-\/-\/-\/-\/-\/-\/-\/-\/-\/-\/-\/-\/-\/-\/-\/-\/-\/-\/-\/-\/-

The 1958-\/-1966 expansion (\(\Delta\ln r = +0.336\)) represents the
opposite condition: all four channels contributing positively, no
offsetting tendencies, and the demand channel leading at 47.7 percent.
Understanding why the demand channel led -\/-\/- and why all other
channels aligned behind it -\/-\/- requires connecting the accounting to
the institutional history. The accumulation dynamics of this period were
shaped by Vietnam-era military expenditure, which sustained aggregate
demand at a scale sufficient to activate productive capacities and drive
nearly half the profit rate recovery \textbackslash citep\{Baker1996\}.
The demand activation operated not merely through the \(\hat{\mu}\)
channel in accounting terms but through its structural consequence: when
productive capacities are more fully utilized, real capital productivity
at normal capacity rises mechanically, because the same capital stock is
generating output closer to its frontier performance -\/-\/- which is
why \(B_{r}\) contributed 22.4 percent of the expansion even without any
acceleration in the underlying rate of technical change. The
still-favourable output-to-capital price environment (the relative price
channel contributing 13.2 percent) amplified the demand effect through
the nominal profit rate. And the distributional channel contributed the
least of the four (16.7 percent), consistent with the
productivity-indexed wage compromise that the Fordist institutional
settlement had routinized: real wages rose with productivity without
improving the profit share, so distribution remained largely inert
during an expansion in which the other three channels were collectively
doing the work. The 1958-\/-1966 period is thus the Fordist accumulation
regime at its most coherent: the institutional settlement holding, the
demand channel leading, and the structural channels responding
proportionately.

\%
-\/-\/-\/-\/-\/-\/-\/-\/-\/-\/-\/-\/-\/-\/-\/-\/-\/-\/-\/-\/-\/-\/-\/-\/-\/-\/-\/-\/-\/-\/-\/-\/-\/-\/-\/-\/-\/-\/-\/-\/-\/-\/-\/-\/-\/-\/-\/-\/-\/-\/-\/-\/-\/-\/-\/-\/-\/-\/-\/-\/-\/-\/-\/-\/-\/-\/-\/-\/-\/-\/-
\% MOVEMENT C3 --- STRUCTURAL CRISIS ONSET 1966-1970 \%
-\/-\/-\/-\/-\/-\/-\/-\/-\/-\/-\/-\/-\/-\/-\/-\/-\/-\/-\/-\/-\/-\/-\/-\/-\/-\/-\/-\/-\/-\/-\/-\/-\/-\/-\/-\/-\/-\/-\/-\/-\/-\/-\/-\/-\/-\/-\/-\/-\/-\/-\/-\/-\/-\/-\/-\/-\/-\/-\/-\/-\/-\/-\/-\/-\/-\/-\/-\/-\/-\/-

The 1966-\/-1970 contraction is where the
decomposition\textquotesingle s analytical architecture pays its fullest
dividend, because what the channel shares record is not merely a large
profit rate decline but the specific structural signature that
distinguishes dysfunctionality from stagnation. No single channel
dominates: the profit share accounts for 34.5 percent of the compounded
decline (\(\Delta\ln r = -0.269\)), real capital productivity for 29.2
percent, and the demand channel for 29.8 percent, with the relative
price channel contributing a minor 6.5 percent. All four channels drag
simultaneously, and no offsetting force is operative. To understand why
this distribution of channel shares constitutes a qualitative break
rather than merely a larger version of the earlier contractions, it is
necessary to hold the contrast with 1951-\/-1954 and 1955-\/-1958 firmly
in view. Those episodes registered one-channel dominance of between 48
and 60 percent precisely because the Fordist institutional settlement
was doing its job: containing distributional conflict through
productivity-indexed wages, sustaining demand through fiscal management,
and holding the organizational conditions of Taylorist production
sufficiently stable to generate real capital productivity growth. When
one channel came under pressure, the others continued to support the
profit rate, and partial reconfiguration could restore the dominant
channel without touching the others. The 1966-\/-1970 contraction
presents a different structure because the institutional containment
failed simultaneously across all three structural channels, leaving no
channel free to offset while the others adjusted. Three mechanisms
converged to produce this outcome. Full-employment conditions sustained
through military expenditure produced labour-market tightening that gave
the working class bargaining leverage it had not possessed since the
pre-Fordist period, directly eroding the profit share and connecting the
distributional channel to the institutional ceiling of full employment
itself \textbackslash citep\{Baker1996, Vidal2019\}. The Taylorist work
organisation that had anchored real capital productivity growth was
reaching its productive limits: the systematic separation of conception
from execution that secured managerial control over the labour process
simultaneously eroded the cooperative and cognitive capacities on which
further productivity improvement depended, as the same deskilling that
had served accumulation in the extensive phase became a fetter on
accumulation in the intensive phase \textbackslash citep\{Braverman1998,
Vidal2022\}. And the demand channel weakened as the Vietnam stimulus
reached its inflationary ceiling and the fiscal space for continued
expansion contracted. The near-equal weight of these three mechanisms in
the decomposition -\/-\/- each accounting for between 29 and 35 percent
of the profit rate contraction -\/-\/- is the empirical expression of
what \textbackslash citet\{Vidal2014\} means by the dysfunctionality of
the accumulation regime: not a failure of any single institutional
arrangement, but the exhaustion of the settlement as a whole,
manifesting as simultaneous pressure across every channel the profit
rate depends on.

\%
-\/-\/-\/-\/-\/-\/-\/-\/-\/-\/-\/-\/-\/-\/-\/-\/-\/-\/-\/-\/-\/-\/-\/-\/-\/-\/-\/-\/-\/-\/-\/-\/-\/-\/-\/-\/-\/-\/-\/-\/-\/-\/-\/-\/-\/-\/-\/-\/-\/-\/-\/-\/-\/-\/-\/-\/-\/-\/-\/-\/-\/-\/-\/-\/-\/-\/-\/-\/-\/-\/-
\% MOVEMENT C4 --- PARTIAL RECOVERY AND OIL SHOCK REATTRIBUTION \%
-\/-\/-\/-\/-\/-\/-\/-\/-\/-\/-\/-\/-\/-\/-\/-\/-\/-\/-\/-\/-\/-\/-\/-\/-\/-\/-\/-\/-\/-\/-\/-\/-\/-\/-\/-\/-\/-\/-\/-\/-\/-\/-\/-\/-\/-\/-\/-\/-\/-\/-\/-\/-\/-\/-\/-\/-\/-\/-\/-\/-\/-\/-\/-\/-\/-\/-\/-\/-\/-\/-

The 1970-\/-1972 expansion (\(\Delta\ln r = +0.066\)) that followed is
the smallest positive swing in the sample and provides a precise reading
of what partial crisis looks like in channel terms. The profit share
leads at 52.7 percent -\/-\/- the distributional adjustment that the
1966-\/-1970 contraction had compressed was partially reversed as
labour-market pressure eased -\/-\/- and capacity utilization
contributes 34.1 percent, with real capital productivity at 23.5
percent. What limits the recovery is visible in the one offsetting
channel: the relative price of output to capital goods contributed
\(-10.4\) percent, meaning that the secular compression of \(Py/PK\) was
already working against profitability recovery even in a period of
distributional adjustment and demand activation. The circuit was
partially restored, but the structural conditions for a renewed Fordist
expansion -\/-\/- a sustained demand environment, an organizational
basis for productivity growth, a favourable price setting -\/-\/- were
not simultaneously in place. This is the partial crisis in
Vidal\textquotesingle s (2014) sense: resolvable in one channel while
the others remain under structural pressure. The 1972-\/-1974
contraction (\(\Delta\ln r =
-0.164\)) delivers the analytically most consequential result of the
four-channel decomposition. In the three-channel nominal framework, this
contraction appears technology-led: the nominal capital productivity
channel accounts for the dominant share of the profit rate decline,
which has historically been interpreted as evidence of a structural
deterioration in the productive efficiency of the US capital stock in
the early 1970s. The four-channel decomposition disaggregates that
finding and records a different structure: the relative price channel
accounts for 57.8 percent of the compounded decline, while real capital
productivity is mildly positive, contributing an offsetting \(-13.0\)
percent of the total swing. What this means is that real productive
efficiency -\/-\/- the capacities of each unit of real capital to
generate real output at normal utilization -\/-\/- continued to grow
during the oil shock contraction. The nominal profit rate was compressed
not because the capital stock became less productive in real terms but
because the price of capital goods rose relative to the price of output,
reducing the nominal rate of return on investment even as the real
productive frontier continued to expand. The oil shock operated through
the \(Py/PK\) channel, not through \(B_{r}\), and its consequence
belongs in the realization crisis cell of the Basu (2019) typology
-\/-\/- specifically, a price-mediated realization crisis rather than a
demand-quantity collapse or a structural technology failure.

\%
-\/-\/-\/-\/-\/-\/-\/-\/-\/-\/-\/-\/-\/-\/-\/-\/-\/-\/-\/-\/-\/-\/-\/-\/-\/-\/-\/-\/-\/-\/-\/-\/-\/-\/-\/-\/-\/-\/-\/-\/-\/-\/-\/-\/-\/-\/-\/-\/-\/-\/-\/-\/-\/-\/-\/-\/-\/-\/-\/-\/-\/-\/-\/-\/-\/-\/-\/-\/-\/-\/-
\% MOVEMENT D1 --- THE CHANNEL PROGRESSION \%
-\/-\/-\/-\/-\/-\/-\/-\/-\/-\/-\/-\/-\/-\/-\/-\/-\/-\/-\/-\/-\/-\/-\/-\/-\/-\/-\/-\/-\/-\/-\/-\/-\/-\/-\/-\/-\/-\/-\/-\/-\/-\/-\/-\/-\/-\/-\/-\/-\/-\/-\/-\/-\/-\/-\/-\/-\/-\/-\/-\/-\/-\/-\/-\/-\/-\/-\/-\/-\/-\/-

Taken together, the four contractions document a structural progression
whose pattern is the analytical core of the present section. The
1951-\/-1954 contraction is \(\pi\)-led (48.5 percent); the 1955-\/-1958
contraction is \(\mu\)-led (60.2 percent); the 1966-\/-1970 contraction
has no dominant channel (\(\pi\) at 34.5, \(B_{r}\) at 29.2, \(\mu\) at
29.8); and the 1972-\/-1974 contraction is \(Py/PK\)-led (57.8 percent).
This progression maps precisely onto Vidal\textquotesingle s (2014)
taxonomy of accumulation regime pathology: single-channel contractions
are the empirical form of stagnation tendencies in
Vidal\textquotesingle s (2014) sense that a functional regime absorbs;
three-channel reinforcing contractions are the empirical form of
structural crisis that no partial reconfiguration can address; and the
price-shock contraction that terminates the sequence is an external
disturbance that accelerated an institutional collapse already underway
rather than initiating it. The capacity to distinguish between these
three categories -\/-\/- and to assign the 1972-\/-1974 contraction to
the third rather than the second -\/-\/- is precisely what the
four-channel decomposition enables and what the three-channel framework
suppresses by attributing the price-mediated contraction to the
technology channel. Figure\textasciitilde\textbackslash ref\{fig:B3\}
makes this progression visible through the normalized channel shares
across all seven swings: the shift from colour concentration in the
early contractions to near-equal distribution across three channels in
the 1966-\/-1970 row, and the subsequent inversion in which green
(\(B_{r}\)) becomes the sole offsetting bar in 1972-\/-1974, captures
the structural transition in a single figure without requiring the
reader to compute the sub-period averages or interpret the swing shares
numerically.

\%
-\/-\/-\/-\/-\/-\/-\/-\/-\/-\/-\/-\/-\/-\/-\/-\/-\/-\/-\/-\/-\/-\/-\/-\/-\/-\/-\/-\/-\/-\/-\/-\/-\/-\/-\/-\/-\/-\/-\/-\/-\/-\/-\/-\/-\/-\/-\/-\/-\/-\/-\/-\/-\/-\/-\/-\/-\/-\/-\/-\/-\/-\/-\/-\/-\/-\/-\/-\/-\/-\/-
\% MOVEMENT D2 --- THE LIMIT OF THE COMPENSATION MECHANISM \%
-\/-\/-\/-\/-\/-\/-\/-\/-\/-\/-\/-\/-\/-\/-\/-\/-\/-\/-\/-\/-\/-\/-\/-\/-\/-\/-\/-\/-\/-\/-\/-\/-\/-\/-\/-\/-\/-\/-\/-\/-\/-\/-\/-\/-\/-\/-\/-\/-\/-\/-\/-\/-\/-\/-\/-\/-\/-\/-\/-\/-\/-\/-\/-\/-\/-\/-\/-\/-\/-\/-

What the sub-period averages and the swing-level decomposition together
establish is not merely a historical narrative of Fordist profitability
dynamics but a precise account of a structural contradiction: the
mechanism through which the Fordist accumulation regime reproduced
profitability contained within it the conditions that brought that
reproduction to a close. The regime sustained the profit rate at
approximately 0.187-\/-0.189 across the Early and Late Fordist
sub-periods through real capital productivity growth -\/-\/- \(B_{r}\)
rising from 0.457 in the Pre-Fordist period to 0.635 in the Late Fordist
-\/-\/- against simultaneous declines in both the activation of
productive capacities (mean \(\hat{\mu}\) falling from 0.875 to 0.723)
and the output-to-capital relative price (mean \(Py/PK\) falling from
1.369 to 1.140). The mechanism was constitutive of the regime rather
than incidental to it: it reflects the overaccumulation character of the
US non-financial corporate sector throughout the Fordist period, with
the transformation elasticity below unity converting each wave of
capital accumulation into productive capacities expansion at a rate that
sustained real capital productivity growth precisely because the
frontier expanded more slowly than the stock. The institutional ceiling
of this mechanism is what the 1966-\/-1970 three-channel contraction
registers. The demand activation that allowed \(B_{r}\) to rise -\/-\/-
because more fully utilized productive capacities generate higher
capital productivity mechanically, not merely through technical change
-\/-\/- was simultaneously generating the full-employment conditions
that eroded \(\pi\) through labour-market tightening and the Taylorist
organizational pressures that exhausted \(B_{r}\) through the
degradation of cooperative productive capacities. The mechanism could
not be sustained precisely because the conditions for its operation were
also the conditions for its dissolution. Demand conditions intervened on
the political economy of the ceiling -\/-\/- the class struggle over
distribution, utilisation, and the institutional forms that govern
whether accumulation can proceed.

\%
-\/-\/-\/-\/-\/-\/-\/-\/-\/-\/-\/-\/-\/-\/-\/-\/-\/-\/-\/-\/-\/-\/-\/-\/-\/-\/-\/-\/-\/-\/-\/-\/-\/-\/-\/-\/-\/-\/-\/-\/-\/-\/-\/-\/-\/-\/-\/-\/-\/-\/-\/-\/-\/-\/-\/-\/-\/-\/-\/-\/-\/-\/-\/-\/-\/-\/-\/-\/-\/-\/-
\% TRANSITION TO §2.7 \%
-\/-\/-\/-\/-\/-\/-\/-\/-\/-\/-\/-\/-\/-\/-\/-\/-\/-\/-\/-\/-\/-\/-\/-\/-\/-\/-\/-\/-\/-\/-\/-\/-\/-\/-\/-\/-\/-\/-\/-\/-\/-\/-\/-\/-\/-\/-\/-\/-\/-\/-\/-\/-\/-\/-\/-\/-\/-\/-\/-\/-\/-\/-\/-\/-\/-\/-\/-\/-\/-\/-

The profitability analysis has established what channels drove the
Fordist profit rate across sub-periods and swings, and has identified
the 1966-\/-1970 three-channel contraction as the structural crisis that
terminated the Fordist compensation mechanism. What it cannot establish
is whether capitalists responded to these channels as independent
signals when deciding how much surplus to reinvest. The decomposition
reveals the accounting structure of profitability;
Stage\textasciitilde C asks the behavioural question that the accounting
structure motivates -\/-\/- whether the demand channel exerted an
independent effect on the recapitalization rate above and beyond its
role through the profit rate, or whether the Okishio restriction
\(\beta_{\mu} = \beta_{B_{r}}\) holds, compressing demand and technology
into a single behavioural lever.

\begin{center}\rule{0.5\linewidth}{0.5pt}\end{center}

\subsection{OUTPUT FORMAT}\label{output-format}

Produce the complete revised LaTeX section, from
\texttt{\textbackslash{}subsubsection\{...\}} to the final transition
sentence, followed by the \texttt{\%\ AUDIT:} block. Nothing else. Do
not explain your choices before or after the LaTeX output.
